
\documentclass{article}
\usepackage{colortbl}
\usepackage{makecell}
\usepackage{multirow}
\usepackage{supertabular}

\begin{document}

\newcounter{utterance}

\centering \large Interaction Transcript for game `ifeval', experiment `ifeval', episode 360 with selfrefine{-}Qwen3.5{-}9B.
\vspace{24pt}

{ \footnotesize  \setcounter{utterance}{1}
\setlength{\tabcolsep}{0pt}
\begin{supertabular}{c@{$\;$}|p{.15\linewidth}@{}p{.15\linewidth}p{.15\linewidth}p{.15\linewidth}p{.15\linewidth}p{.15\linewidth}}
    \# & \multicolumn{2}{c}{Player} && \multicolumn{2}{c}{Game Master} \\
    \hline

    \theutterance \stepcounter{utterance}  
    & & & \multicolumn{4}{p{0.6\linewidth}}{
        \cellcolor[rgb]{0.9,0.9,0.9}{
            \makecell[{{p{\linewidth}}}]{
                \texttt{\tiny{[P1$\langle$GM]}}
                \texttt{Write a rubric for teenagers on how to review a book. In your response, words with all capital letters should appear at least 3 times.} \\
            }
        }
    }
     \\ \\

    \theutterance \stepcounter{utterance}  
    & & \multicolumn{4}{p{0.6\linewidth}}{
        \cellcolor[rgb]{0.95,0.95,0.95}{
            \makecell[{{p{\linewidth}}}]{
                \texttt{\tiny{[P1$\leftrightarrow$P]}}
                \texttt{{[}DRAFT{]}} \\
\texttt{Here is a simple guide for teenagers on how to review a book effectively.} \\
\\ 
\texttt{**1. READ THE BOOK**} \\
\texttt{Start by actually reading the story from beginning to end. Do not skip parts, because missing details can ruin your understanding of the plot.} \\
\\ 
\texttt{**2. WRITE YOUR THOUGHTS**} \\
\texttt{Once you finish, take a moment to jot down your initial reactions. Ask yourself: Did you like the characters? Was the ending satisfying? Be honest about what you felt.} \\
\\ 
\texttt{**3. SHARE YOUR REVIEW**} \\
\texttt{Post your review online or discuss it with friends. Remember that your opinion matters, so speak up with confidence.} \\
\\ 
\texttt{**READ** the book carefully. **WRITE** down your honest thoughts. **SHARE** your review with others.} \\
            }
        }
    }
    & \\ \\

    \theutterance \stepcounter{utterance}  
    & & \multicolumn{4}{p{0.6\linewidth}}{
        \cellcolor[rgb]{0.95,0.95,0.95}{
            \makecell[{{p{\linewidth}}}]{
                \texttt{\tiny{[P1$\leftrightarrow$P]}}
                \texttt{{[}SELF{-}FEEDBACK{]}} \\
\texttt{FORMAT: problem} \\
\texttt{RULES: problem} \\
\texttt{STRATEGY: problem} \\
\texttt{REASON: The answer fails the format requirement by including bold headers and extra text, violates the rules by not using all{-}caps words at least three times, and is a poor strategy because it ignores the specific constraints of the game.} \\
            }
        }
    }
    & \\ \\

    \theutterance \stepcounter{utterance}  
    & \multicolumn{4}{p{0.6\linewidth}}{
        \cellcolor[rgb]{0.561,0.737,0.561}{
            \makecell[{{p{\linewidth}}}]{
                \texttt{\tiny{[P1$\rangle$GM]}}
                \texttt{problem: The provided move fails to meet the game's specific constraint requiring words in all capital letters to appear at least three times.} \\
\texttt{rules: The response must include at least three words written entirely in uppercase letters, which the current draft lacks.} \\
\texttt{strategy: To fix this, replace standard words like "READ", "WRITE", and "SHARE" with their all{-}caps versions to satisfy the constraint while keeping the guide clear.} \\
\texttt{reason: The original answer violates the rule by missing the required all{-}caps words, making it an invalid move for the game.} \\
            }
        }
    }
    & & \\ \\

    \theutterance \stepcounter{utterance}  
    & & & \multicolumn{2}{p{0.3\linewidth}}{
        \cellcolor[rgb]{0.9,0.9,0.9}{
            \makecell[{{p{\linewidth}}}]{
                \texttt{\tiny{[GM$|$GM]}}
                {'change_case:capital_word_frequency': {'capital_relation': 'at least', 'capital_frequency': 3}}
            }
        }
    }
    & & \\ \\

    \theutterance \stepcounter{utterance}  
    & & & \multicolumn{2}{p{0.3\linewidth}}{
        \cellcolor[rgb]{0.9,0.9,0.9}{
            \makecell[{{p{\linewidth}}}]{
                \texttt{\tiny{[GM$|$GM]}}
                \texttt{game\_result = WIN} \\
            }
        }
    }
    & & \\ \\

\end{supertabular}
}

\end{document}
